\section{Limitations and Responsible Interpretation}
\label{sec:limitations}

\begin{limitbox}[What the paper does not establish]
\begin{itemize}
  \item It does not establish novelty of coherent measurement deferral or
  teleportation cleanup; those follow from known principles.
  \item It does not establish an attack on BB84 under the standard channel-only
  threat model. The fixed-input result assumes access to source preparation
  controls or an equivalent privileged implementation boundary.
  \item It does not yet identify the exact Boolean acceptance register of the
  complete original educational detector. The current acceptance rule is
  provisional.
  \item It does not demonstrate malware on production quantum hardware, and the
  validated reference backend is an ideal NumPy model rather than a device run.
  \item It does not yet provide full branch-level Qiskit parity or a malicious
  compiler prototype.
  \item It does not show that encrypted data can always be decrypted or usefully
  exploited after copying.
  \item Agreement between the coherent cleanup implementation and the SWAP
  reference on $\ket{\psi00}$ does not imply global unitary equivalence on
  arbitrary three-qubit inputs.
\end{itemize}
\end{limitbox}

The paper instead offers a carefully bounded synthesis: a provable cleanup
identity, an exact export-level state comparison, a fixed-input source-access
leakage result, a general information-disturbance boundary, and a continuously
tested research program for secure qubit reuse.
